% §19.6 The Classical Sector of a Quantum Formalism
%
% Monograph register: make the QM connection precise and honest. The formalism
% here is identical to finite-dimensional QM EXCEPT the vectors are probabilities
% (or 0/1 indicators), not amplitudes; no Born-rule square, no measurement
% collapse. The Bernoulli framework is exactly the classical stochastic-matrix
% algebra on the DIAGONAL of the density operator, which is all a measurement
% with a diagonal POVM sees, for ANY state: Tr(M_b rho) = sum_t A(b|t) rho_tt
% (scope corrected 2026-10; earlier drafts restricted it to diagonal states). Bra-ket and tensor notation
% are finite-dim linear algebra, not QM-specific. Tie: sec:quantum-contextuality
% (ch16) is where the diagonal-sector representation stops being faithful (need
% amplitudes / off-diagonal coherence); rem:quantum-amplification (ch11) is where
% the square reappears. State plainly: adopting bra-ket here is a notational
% convenience for the classical algebra, not a quantum claim. Cross-refs
% sec:n-ary-channels, def:n-ary-channel, sec:quantum-contextuality,
% rem:quantum-amplification, sec:operator-computational-basis,
% def:set-noise-operator, prop:parity-spectral, sec:operator-what-for. DEFINED:
% sec:operator-classical-sector.

\section{The Classical Sector of a Quantum Formalism}
\label{sec:operator-classical-sector}

The notation of this chapter, kets, tensor products, eigenvalues, is the
notation of finite-dimensional quantum mechanics, and the resemblance is close
enough that the relationship deserves a precise statement. The framework is the
classical sector\index{classical sector} of the quantum formalism: the same linear algebra, on
probabilities rather than amplitudes, acting on the diagonal of the state,
which is everything a computational-basis measurement sees.
This section says exactly what coincides, what does not, and why the bra-ket
notation is a convenience and not a claim.

\subsection*{Same algebra, different objects}

Set the two formalisms side by side. Finite-dimensional quantum mechanics works
in $\ComplexSet^{d}$, builds composite systems by the tensor product
$\ComplexSet^{d_1} \otimes \ComplexSet^{d_2}$, evolves states by linear
operators, and reads spectra off those operators. The operator view of this
chapter does all four: it works in $\ComplexSet^{N}$ for maps and
$(\ComplexSet^{2})^{\otimes N}$ for sets
(\Cref{sec:operator-computational-basis}, \Cref{def:set-noise-operator}),
composes systems by the Kronecker product, evolves by stochastic matrices, and
reads error rates off eigenvalues (\Cref{prop:parity-spectral}). The
linear-algebraic machinery is identical.

The objects acted on are not. A quantum state vector holds probability
\emph{amplitudes}, complex numbers whose squared magnitudes are probabilities,
and amplitudes can be negative or complex and so can cancel. The vectors here
hold probabilities directly (a simplex vector for a random element) or bits (a
$0/1$ occupation vector for a set), and probabilities only accumulate. Two
consequences of quantum mechanics are therefore simply absent. There is no
Born-rule square: a coordinate of a vector here \emph{is} a probability, not an
amplitude to be squared into one. And there is no measurement collapse: reading
a coordinate does not disturb the others, because the vector is already a
classical description, not a superposition awaiting measurement. The chapter
uses kets to write one-hot indicators and tensor products to write independence,
and that is the whole of the borrowing.

\subsection*{The framework lives on the diagonal}

The precise locating statement uses the density operator. A quantum
system in $\ComplexSet^{d}$ is described in general by a density matrix
$\rho$\index{density matrix}, a positive operator of unit trace whose
diagonal entries $\rho_{tt}$ are the probabilities of the basis outcomes
and whose off-diagonal entries $\rho_{ts}$ are the coherences that encode
interference. A measurement is described by a POVM $\{M_b\}$, positive
operators summing to the identity, with outcome probabilities
$\operatorname{Tr}(M_b \rho)$ \parencite{nielsen-chuang-2010}. When the
detector's POVM elements are diagonal in the computational basis\index{detector POVM},
\[
  M_b \;=\; \sum_t A(b \mid t)\, \ketbra{t}{t},
  \qquad
  \operatorname{Tr}(M_b \rho) \;=\; \sum_t A(b \mid t)\, \rho_{tt}
  \quad \text{for every } \rho,
\]
so the outcome law is the column-stochastic assignment matrix $A$ applied
to the diagonal of $\rho$, and the coherences never enter. This holds for
an eigenstate, a mixture, or a coherent superposition alike; it is a fact
about the measurement, not about the state. A probability vector here is
the diagonal of a density matrix, and the stochastic matrix is the
assignment matrix of a diagonal measurement. This is the readout instance
of \Cref{sec:n-ary-channels}\index{quantum readout}: measuring $n$ qubits in
the computational basis is an n-ary channel (\Cref{def:n-ary-channel})
whose assignment matrix acts on diagonals, and when per-qubit readout
errors are independent it is the Kronecker product of single-qubit
channels, the same factorization \Cref{def:set-noise-operator} gives a set.
The readout-error literature uses exactly these objects: tensor-product
and correlated assignment matrices, their inversion, and their
calibration
\parencite{bravyi-et-al-2021,nation-et-al-2021,smith-et-al-2021,maciejewski-zimboras-oszmaniec-2020,maciejewski-et-al-2021,hamilton-et-al-2020}.
For the outcome statistics of diagonal measurements the framework is
exact, whatever the state.

\subsection*{Where the diagonal description stops being faithful}

What the diagonal cannot answer is any question whose answer depends on
the coherences, and the boundary is the subject of
\Cref{sec:quantum-contextuality}. Two kinds of question cross it. One asks
about a measurement that is not diagonal in the basis defining the latent
value: a rotated basis, an entangling measurement, or a coherent
operation applied before readout. There the outcome law depends on the
off-diagonal entries of $\rho$, and the diagonal alone does not determine
it. The other asks for joint statistics of incompatible measurements, and
there may be no consistent assignment of definite values to all the
observables at once to which a stochastic matrix could be applied: the
regime where no latent value exists, which
\Cref{sec:quantum-contextuality} shows, through Fine's theorem and the Bell
and Kochen-Specker obstructions, to be a falsifiable physical fact and not
a modeling convenience. In operator terms the framework's vectors are
diagonals, a diagonal measurement sees nothing else, and contextuality
lives in what a single diagonal cannot represent. No decoherence is needed
for the framework to apply to a diagonal measurement; what decoherence
removes are the coherences that \emph{other} measurements could have
detected.

The other half of the boundary is where the square returns.
\Cref{rem:quantum-amplification} compared classical amplification\index{error amplification}, which composes
probabilities and reaches near-certainty in $O(1/p)$ repeated looks, with
amplitude amplification\index{amplitude amplification}, which composes amplitudes and reaches it in
$O(1/\sqrt{p})$ steps. The quadratic gap is the Born-rule square applied to a
quantity that accumulates linearly, and it appears because amplitude
amplification acts coherently between preparation and readout, on the
amplitude-bearing part of the state that the diagonal description
discards. The framework's amplification theorems are the
diagonal-sector statements; the quantum speedup is what the off-diagonal sector
buys, and it is unavailable to a formalism that keeps only the diagonal.

\subsection*{The notation is a convenience}

The honest summary is therefore narrow. Adopting bra-ket and tensor notation in
this chapter is a notational convenience for the classical stochastic algebra:
it makes serial composition a product (\Cref{prop:operator-map-composition}),
independence a factorization (\Cref{prop:set-operator-is-axiom-1}), and a parity
error rate a product of eigenvalues (\Cref{prop:parity-spectral}), all of which
are awkward to write in the predicate language and natural in the operator one.
It is not a claim that the framework is quantum mechanical. The framework is the
classical, diagonal, real-and-nonnegative part of the quantum formalism;
\Cref{sec:operator-what-for} opened by calling the operator view a second lens
on the same set-theoretic content, and the precise version of that claim is this:
the lens is finite-dimensional linear algebra, the same linear algebra quantum
mechanics uses, pointed at the diagonal where probabilities live and amplitudes
do not.
